\documentclass[10pt]{article}
\usepackage[margin=1in]{geometry}
\usepackage{amsmath}
\usepackage{amssymb}
\usepackage{bm}
\usepackage{bbm}
\usepackage{booktabs}
\usepackage{graphicx}
\usepackage{subcaption}
\usepackage{float}
\usepackage{hyperref}
\usepackage{enumitem}

\title{Bid-Share GMM and Cost Recovery Audit}
\author{Konan Hara}
\date{September 13, 2026}

\begin{document}
\maketitle

\section{Purpose and Current Status}

This note records the preferred first-stage GMM estimator and the aligned
second-stage cost-recovery procedure. Both stages have been reproduced on the
HPCC. The reported standard errors and central 95 percent percentile intervals
use the completed 200-draw global firm-pairs bootstrap, conditional on the
observed auctions.

The Step 1 sample contains 741 retained winner observations from 22 auctions.
Auctions 23 and 24, the two 2021 auctions conducted during the COVID-period
market disruption, are excluded upstream. The sample contains 600 winners from
16 sealed-bid auctions and 141 winners from six descending-clock auctions.
Every retained winner receives equal weight. Auction-recorded capacity defines
the bid share and its bounds; SIGA-adjusted capacity is not used.

\section{Prior Started Wind Capacity}

Let $K^{W}_{fa}$ denote the total physical assurance of wind projects
previously awarded to firm $f$ whose contractual delivery periods had begun by
auction $a$. Each project is counted once at its first award. The construction
uses only public auction-award information available before auction $a$:
\[
K^{W}_{fa}
=
\sum_{j:\,f(j)=f}
C_j^A
\mathbbm{1}\{a(j)<a,\ t^{\mathrm{start}}_j\leq t_a\}.
\]
The level $K^W$ is measured in MW. Its mean is 54.9 MW, its median is zero,
and its range is 0--421.7 MW in the estimation sample. The preferred regressor
is
\begin{equation}
K^{\log}_{fa}=\log\left(1+\frac{K^W_{fa}}{100}\right).
\label{eq:klog}
\end{equation}
This transformation preserves the structural zero for firms with no prior
started wind capacity and allows diminishing marginal effects at larger
histories.

Effective absolute risk aversion is parameterized as
\begin{equation}
\gamma_i
=
\exp\left(
\alpha_{\gamma0}
+\alpha_{\gamma K}K^{\log}_{f(i),a(i)}
\right).
\label{eq:gamma}
\end{equation}
Expected wholesale-market profitability is
\begin{equation}
\mu_{r,i}=\alpha_{r0}+\alpha_{r1}m_i,
\label{eq:mur}
\end{equation}
where $m_i$ is forecasted marginal cost of electricity.

\section{Reduced-Form Evidence}

Table \ref{tab:reduced_form} reports OLS regressions of the bid share on
purchase-agreement price and forecasted marginal cost. Standard errors are
two-way clustered by auction and firm. For transparency, Column (2) adds prior
started wind capacity in levels, $K^W_{fa}/100$. The structural model continues
to use the transformed regressor in equation~\eqref{eq:klog}.

\begin{table}[H]
\centering
\caption{Bid share}
\label{tab:reduced_form}
\begin{tabular}{lrr}
\toprule
& (1) & (2)\\
\midrule
Purchase-agreement price
& 0.01264 & 0.01161\\
& (0.00418) & (0.00387)\\
Forecasted marginal cost
& -0.01762 & -0.01774\\
& (0.00717) & (0.00685)\\
Prior started wind capacity / 100 MW
& & -0.03288\\
& & (0.01995)\\
\midrule
Observations & 741 & 741\\
Auctions & 22 & 22\\
Firms & 63 & 63\\
$R^2$ & 0.400 & 0.420\\
\bottomrule
\end{tabular}
\end{table}

The bid-price coefficient remains positive and statistically significant at
the one-percent level ($p=0.0068$) after including the history measure. The
history coefficient is negative but does not reach the ten-percent level under
two-way clustering ($p=0.114$). Its sign nevertheless
provides suggestive reduced-form motivation for allowing firms with larger
prior wind portfolios to be more risk tolerant.

\section{GMM Estimator}

For each bidder the estimator uses three conditional moment blocks: the lower
bound probability, the upper bound probability, and the conditional mean among
interior bid shares. It does not include a separate variance moment and does
not apply participant-count weights.

The estimator proceeds through three nested models:
\begin{enumerate}[label=(\roman*)]
\item homogeneous $\gamma$ and homogeneous $\mu_r$;
\item homogeneous $\gamma$ and $\mu_{r,i}$ as in \eqref{eq:mur};
\item heterogeneous $\gamma_i$ as in \eqref{eq:gamma} and $\mu_{r,i}$ as in
\eqref{eq:mur}.
\end{enumerate}

All original-sample optimizations use bounded Nelder--Mead as implemented in
NLopt. The first model begins with an identity weighting matrix and is
initialized at $(\gamma,\mu_r,\sigma_\eta)=(0.5,25,0.1)$. Each subsequent
model begins from the preceding model's second-step estimates and a weighting
matrix evaluated at those estimates. Within every original-sample model, the
second GMM step starts from the first-step estimate and uses the moment
covariance matrix evaluated at that estimate.

The restricted model uses the four instruments
\[
\mathbbm{1}\{SB\},\quad
\mathbbm{1}\{DC\},\quad
b_i\mathbbm{1}\{SB\},\quad
b_i\mathbbm{1}\{DC\}.
\]
The homogeneous-$\gamma$ model adds $m_i$. The heterogeneous model adds the
single pooled interaction
\[
b_i K^{\log}_{f(i),a(i)}.
\]
Applying each instrument to all three moment blocks gives 12, 15, and 18
moments, respectively.

\section{Structural Estimates}

\begin{table}[H]
\centering
\caption{First-stage GMM estimates}
\label{tab:gmm_estimates}
\begin{tabular}{lrrrr}
\toprule
& \multicolumn{2}{c}{Homogeneous $\gamma$}
& \multicolumn{2}{c}{Heterogeneous $\gamma$}\\
\cmidrule(lr){2-3}\cmidrule(lr){4-5}
& Step 1 & Step 2 & Step 1 & Step 2\\
\midrule
$\alpha_{\gamma0}$ & -0.7671 & -0.7083 & -0.5810 & -0.5530\\
$\alpha_{\gamma K}$ & & & -0.1165 & -0.1205\\
$\alpha_{r0}$ & 2.2906 & 2.5505 & 1.2814 & 1.0634\\
$\alpha_{r1}$ & 0.6933 & 0.6814 & 0.7224 & 0.7347\\
$\alpha_\eta$ & -1.4679 & -1.5532 & -1.6470 & -1.7234\\
Objective & 0.3404 & 0.3286 & 0.3665 & 0.3650\\
\bottomrule
\end{tabular}
\end{table}

All six optimizations in the nested sequence report
\texttt{LOCALLY\_SOLVED}, with NLopt status \texttt{FTOL\_REACHED}. The
homogeneous second-step estimate implies $\gamma=0.492$. Under the final
heterogeneous estimate, the winner-sample gamma distribution has mean 0.555,
standard deviation 0.030, median 0.575, 10th percentile 0.512, 90th percentile
0.575, and range 0.471--0.575. The first-step and second-step means are 0.540
and 0.555, respectively. The final estimate also implies
$\widehat\sigma_\eta=0.178$ and an observed-sample $\mu_r$ range of
23.62--35.89.

The point estimate $\widehat\alpha_{\gamma K}=-0.121$ associates a larger
prior wind portfolio with lower risk aversion, while leaving the mean close to
the homogeneous-$\gamma$ estimate. The modest cross-sectional dispersion is
economically attractive for treating heterogeneous risk aversion as an
important extension of the homogeneous model rather than as a different
average-risk specification.

\section{Bootstrap Inference}

The preferred bootstrap conditions on the 22 observed auctions and resamples
the 63 firms globally with replacement. A firm receives one draw-specific
multiplicity that applies to all of its observations across auctions and
auction formats. This preserves within-firm dependence while holding auction
identities and characteristics fixed. The bootstrap moments are recentered at
each model's original-sample second-step estimate.

The original estimator remains the two-step GMM estimator described above. In
the bootstrap, the model-specific second-step weighting matrix is held fixed
across draws. Each matrix is recalculated from the original sample at the
corresponding original first-step estimate. The bootstrap therefore captures
sampling variation in the empirical moments and firm composition without
re-estimating a potentially ill-conditioned optimal weighting matrix in every
draw. Under conventional large-cluster GMM asymptotics, estimation of a
consistent weighting matrix does not affect the estimator's first-order
distribution under correct specification. The fixed-matrix procedure does not
reproduce finite-sample variation in the feasible weighting matrix; this is the
main inferential qualification.

\begin{table}[H]
\centering
\caption{Fixed-weight global firm-pairs bootstrap inference}
\label{tab:gmm_bootstrap_parameters}
\begin{tabular}{lrrrr}
\toprule
& \multicolumn{2}{c}{Homogeneous $\gamma$}
& \multicolumn{2}{c}{Heterogeneous $\gamma$}\\
\cmidrule(lr){2-3}\cmidrule(lr){4-5}
& Estimate & Bootstrap SE & Estimate & Bootstrap SE\\
\midrule
$\alpha_{\gamma0}$ & -0.7083 & 0.0678 & -0.5530 & 0.2630\\
$\alpha_{\gamma K}$ & & & -0.1205 & 0.3369\\
$\alpha_{r0}$ & 2.5505 & 2.9882 & 1.0634 & 3.7464\\
$\alpha_{r1}$ & 0.6814 & 0.0829 & 0.7347 & 0.1069\\
$\alpha_\eta$ & -1.5532 & 0.0836 & -1.7234 & 0.1165\\
\bottomrule
\end{tabular}
\begin{minipage}{0.92\textwidth}
\footnotesize
\textit{Notes:} Standard errors are standard deviations across 200 bootstrap
draws. Firms are resampled globally with replacement; auctions are held fixed;
moments are recentered; and each model uses its original-sample second-step
weighting matrix.
\end{minipage}
\end{table}

\begin{table}[H]
\centering
\caption{Bootstrap inference for implied risk aversion}
\label{tab:gmm_bootstrap_gamma}
\begin{tabular}{lrrr}
\toprule
Model & Estimate & Bootstrap SE & Central 95\% interval\\
\midrule
Homogeneous $\gamma$, homogeneous $\mu_r$
& 0.391 & 0.050 & [0.300, 0.492]\\
Homogeneous $\gamma$
& 0.492 & 0.033 & [0.427, 0.557]\\
Heterogeneous $\gamma$ (sample mean)
& 0.555 & 0.168 & [0.464, 1.028]\\
\bottomrule
\end{tabular}
\begin{minipage}{0.92\textwidth}
\footnotesize
\textit{Notes:} Intervals are the 2.5th and 97.5th percentiles of the 200
bootstrap estimates. The heterogeneous statistic averages $\gamma_i$ over the
original estimation sample in each draw.
\end{minipage}
\end{table}

The central 95 percent percentile interval for $\alpha_{\gamma K}$ is
$[-1.128,0.324]$; 156 of the 200 estimates are negative. The heterogeneity
coefficient is therefore imprecise, though its bootstrap distribution leans
toward the economically expected sign. The heterogeneous-minus-homogeneous
mean-gamma difference is 0.063, with central interval $[-0.015,0.531]$.
The central interval for the heterogeneous gamma standard deviation is
$[0.001,0.395]$, so dispersion is less precisely determined than the mean.

Although $\alpha_{r0}$ has a large marginal standard error, it is strongly
negatively correlated with $\alpha_{r1}$ across draws ($-0.987$). Consequently,
the fitted $\mu_r$ endpoints are much more stable: their central intervals are
approximately [22.10, 24.99] and [32.79, 37.94]. The winner-sample mean full
wholesale-market risk premium is 8.33, with bootstrap standard error 2.58 and
central interval [6.92, 15.61].

All 600 bootstrap optimizations are \texttt{LOCALLY\_SOLVED} and terminate
with \texttt{FTOL\_REACHED}. No estimate is within one percent of a
computational parameter bound.

\subsection{Step 1 Sensitivity Analysis}

All 18 sensitivity specifications and their continuation steps converge. Table
\ref{tab:step1_sensitivity} summarizes the ranges for the final heterogeneous
mean gamma and winner-sample mean full-exposure risk premium.

\begin{table}[H]
\centering
\small
\caption{Step 1 sensitivity ranges}
\label{tab:step1_sensitivity}
\begin{tabular}{lrrr}
\toprule
Perturbation & Grid & Mean $\gamma$ & Risk premium\\
\midrule
Spot-price variance multiplier & 0.5--1.5 & 0.245--2.318 & 8.256--8.690\\
Spot-price mean shift & $-2$--$2$ SD & 0.451--0.612 & 6.769--9.174\\
Winning-bid variance multiplier & 0.5--1.5 & 0.548--0.560 & 8.217--8.405\\
Discount factor $\delta$ & 0.90--0.99 & 0.329--0.882 & 5.252--12.222\\
\bottomrule
\end{tabular}
\end{table}

As in Step 2 below, gamma adjusts sharply to alternative spot-price variances,
but the implied risk premium remains stable. Results are also insensitive to
the winning-bid variance. The discount-factor comparison requires a scale
adjustment because the risk premium is itself a discounted average over the
contract horizon. Table \ref{tab:step1_delta_scaling} divides each winner's
risk premium by its auction-specific
$\widetilde\delta_a=T^{-1}\sum_{t=l_a}^{l_a+T-1}\delta^t$ before averaging.

\begin{table}[H]
\centering
\caption{Step 1 discount-factor sensitivity in comparable units}
\label{tab:step1_delta_scaling}
\begin{tabular}{rrrr}
\toprule
$\delta$ & Mean $\widetilde\delta_a$ & Risk premium
& Risk premium / $\widetilde\delta_a$\\
\midrule
0.90 & 0.309 & 5.252 & 16.754\\
0.95 & 0.539 & 8.328 & 15.371\\
0.99 & 0.880 & 12.222 & 13.888\\
\bottomrule
\end{tabular}
\begin{minipage}{0.92\textwidth}
\footnotesize
\textit{Notes:} Statistics average equally over the 741 Step 1 winners. The
last column first divides each observation's risk premium by its own auction's
$\widetilde\delta_a$ and then averages.
\end{minipage}
\end{table}

The unscaled premium changes by a factor of 2.33 between the endpoints, whereas
the normalized premium changes by a factor of only 1.21. The residual movement
is reasonable: changing $\delta$ also changes the discounted wholesale-price
variance and the re-estimated gamma, so the risk premium need not be exactly
proportional to $\widetilde\delta_a$.

\section{Goodness of Fit}

\begin{table}[H]
\centering
\caption{Fit of the final heterogeneous estimate}
\label{tab:gof}
\begin{tabular}{lrrrr}
\toprule
Sample & Observed mean & Predicted mean & RMSE & Category $L^1$\\
\midrule
All & 0.880 & 0.852 & 0.162 & 0.223\\
Sealed bid & 0.943 & 0.952 & 0.110 & 0.298\\
Descending clock & 0.615 & 0.429 & 0.294 & 0.095\\
\bottomrule
\end{tabular}
\end{table}

The aggregate and sealed-bid means fit closely. The model understates the mean
descending-clock bid share. Its aggregate DC category frequencies fit better:
the observed and predicted lower-bound frequencies are 0.362 and 0.399, while
the observed and predicted interior frequencies are 0.589 and 0.599. The model
nearly eliminates the relatively rare DC upper-bound outcomes, predicting
0.002 against 0.050 observed.

\begin{figure}[H]
\centering
\includegraphics[width=\textwidth]{bidshare_gmm_hist_gof.png}
\caption{Observed and predicted bid-share distributions}
\end{figure}

\begin{figure}[H]
\centering
\includegraphics[width=\textwidth]{bidshare_gmm_auction_mean_gof.png}
\caption{Observed and predicted mean bid shares by auction}
\end{figure}

\begin{figure}[H]
\centering
\includegraphics[width=\textwidth]{bidshare_gmm_dc_auction_bounds_gof.png}
\caption{Descending-clock bound frequencies by auction}
\end{figure}

\section{Step 2 Alignment and Cost Recovery}

Step 2 uses the sealed-bid auctions to estimate equilibrium winning
probabilities, invert the bid-price first-order condition for winner costs, and
simulate the participant distribution of $(\gamma,c)$. It uses the same
economic definitions and bootstrap sampling unit as Step 1.

\subsection{Auction Capacity and Firm History}

Auction-recorded physical assurance $C_i^A$ defines the portfolio share
$q_i=Q_i/C_i^A$, the feasible quantity bounds, and the awarded quantity
$q_iC_i^A$ entering the winning rule. For factual winners, risk aversion is
evaluated directly using equation~\eqref{eq:gamma}.

For participants, the concomitant argument recovers the joint distribution of
prior-started wind capacity and current auction capacity. Let $b^*$ denote a
participant's latent equilibrium bid. Within each auction,
\[
f_{K,C,b^*\mid a}(K,C,b)
=f_{K,C\mid b^*,a}(K,C\mid b)f_{b^*\mid a}(b).
\]
The order-statistic likelihood recovers $f_{b^*\mid a}$ from winner bids and
participant counts. Conditional on $b^*$, winner concomitants identify
$f_{K,C\mid b^*,a}$ because selection changes the marginal bid distribution,
not the conditional distribution of characteristics given the bid.

The implementation uses
\[
f_{K,C\mid b^*,a}=f_{K\mid C,b^*,a}f_{C\mid b^*,a}.
\]
Current auction capacity follows
\[
\log C_i^A
=\beta_{Cb}b_i^*+\beta_{Ca}+\varepsilon_{Ci},
\qquad \varepsilon_{Ci}\sim\mathcal N(0,\sigma_C^2).
\]
For prior-started wind capacity, the censored latent model is
\[
(K^{\log}_i)^*
=\beta_{Kb}b_i^*+\beta_{KC}\log C_i^A+\beta_{Ka}
+\varepsilon_{Ki},
\qquad
K^{\log}_i=\max\{0,(K^{\log}_i)^*\},
\qquad
\varepsilon_{Ki}\sim\mathcal N(0,\sigma_K^2).
\]
A single auction fixed effect governs both the probability mass at $K^W=0$
and the positive-history distribution. The first six sealed-bid auctions occur
before any previously awarded wind contract could have started, so $K^W=0$ is
imposed structurally. Applying
$\gamma=\exp(\alpha_{\gamma0}+\alpha_{\gamma K}K^{\log})$ induces the joint
participant distribution of risk aversion, capacity, and equilibrium bids.

The original-sample concomitant models both converge locally. The estimated
conditional variances are $\widehat\sigma_C^2=0.127$ and
$\widehat\sigma_K^2=0.537$. The common slopes in the prior-capacity equation,
$\widehat\beta_{Kb}=-0.046$ and $\widehat\beta_{KC}=0.036$, are imprecisely
estimated; the model's main role is to preserve the observed mass at zero and
the joint dependence needed for participant simulation, rather than to assign
a structural interpretation to these nuisance coefficients.

\subsection{Winner Cost Recovery}

The two sealed-bid auctions without participant counts are excluded from
winning-probability construction and cost recovery. The factual Step 2 sample
therefore contains 549 winners from 14 sealed-bid auctions, representing 4,776
participants and 119,746.7 MW of participant capacity.

\begin{table}[H]
\centering
\caption{Capacity-weighted winner outcomes}
\label{tab:winner_results}
\begin{tabular}{lrrr}
\toprule
Outcome & Estimate & Bootstrap SE & Central 95\% interval\\
\midrule
Risk aversion $\gamma$ & 0.561 & 0.208 & [0.456, 1.250]\\
Optimal bid share & 0.975 & 0.007 & [0.962, 0.989]\\
Full-exposure risk premium & 8.711 & 3.252 & [7.070, 19.180]\\
Markup & 0.950 & 0.154 & [0.630, 1.264]\\
Cost & 23.987 & 0.587 & [22.952, 25.133]\\
\bottomrule
\end{tabular}
\begin{minipage}{0.92\textwidth}
\footnotesize
\textit{Notes:} Outcomes average over factual winners using awarded capacity.
Intervals are the 2.5th and 97.5th percentiles of the bootstrap distribution.
\end{minipage}
\end{table}

\subsection{Participant Distribution}

The participant distribution is summarized by a joint linear projection of
$(\log\gamma,c)$ on the existing polynomial in auction date, lead time, and
their interactions. With $S_a$ valid simulations in auction $a$, the three
weights are
\[
w^A_{sa}=\frac{1}{S_a},
\qquad
w^N_{sa}=\frac{N_a}{S_a},
\qquad
w^C_{sa}=TC_a\frac{C_{sa}}{\sum_{r=1}^{S_a}C_{ra}},
\]
where $TC_a$ is the observed total capacity of all participants. These
describe the average auction, average participant, and average MW,
respectively.

\begin{table}[H]
\centering
\small
\caption{Recovered participant distribution}
\label{tab:participant_results}
\begin{tabular}{lrrrr}
\toprule
Weighting & Mean & Central 95\% interval & SD & Central 95\% interval\\
\midrule
\multicolumn{5}{l}{\textit{Panel A: Risk aversion $\gamma$}}\\
Equal auction & 0.565 & [0.456, 1.254] & 0.024 & [0.001, 0.274]\\
Participant & 0.564 & [0.462, 1.208] & 0.024 & [0.001, 0.292]\\
Capacity & 0.564 & [0.461, 1.210] & 0.024 & [0.001, 0.283]\\
\addlinespace
\multicolumn{5}{l}{\textit{Panel B: Cost}}\\
Equal auction & 27.007 & [26.225, 27.232] & 5.433 & [5.048, 5.508]\\
Participant & 26.464 & [25.929, 26.758] & 5.364 & [5.051, 5.434]\\
Capacity & 26.618 & [26.038, 27.006] & 5.468 & [5.184, 5.535]\\
\bottomrule
\end{tabular}
\begin{minipage}{0.92\textwidth}
\footnotesize
\textit{Notes:} Cost is discounted to the auction date. Central intervals are
the 2.5th and 97.5th percentiles across 200 bootstrap draws.
\end{minipage}
\end{table}

The three weighting conventions give nearly identical means and standard
deviations. The point estimate of the gamma standard deviation is about 0.024,
although its percentile interval is wide and asymmetric; heterogeneity is much
less precisely estimated than average gamma. The participant-weighted cost
standard deviation is 5.36 $/MWh$.
Undiscounted mean costs are 31.286, 30.884, and 31.053 under equal-auction,
participant, and capacity weights. The estimated correlations between gamma
and cost are small: $-0.027$, 0.079, and 0.096. Their central intervals,
respectively $[-0.275,0.291]$, $[-0.231,0.335]$, and $[-0.225,0.334]$, all
include zero. The average gamma level is therefore much more stable than its
cross-sectional variance or its covariance with cost. All 200 draws deliver
finite estimates for every aggregate summary and every linear-projection
coefficient.

Figures \ref{fig:gamma_trends} and \ref{fig:cost_trends} plot the equal-auction
linear projections, matching the weighting convention used in the paper's
existing cost-trend figure. The gamma projection is linear in $\log\gamma$;
Panel (a) exponentiates the fitted value, while Panel (b) reports the implied
level effect of an additional lead-time year. The confidence bands are
pointwise central 95 percent bootstrap percentile intervals.

\begin{figure}[H]
\centering
\begin{subfigure}{0.49\textwidth}
\centering
\includegraphics[width=\textwidth]{step2_gmm_gamma_t_l0.png}
\caption{Auction-date trend at zero lead time}
\end{subfigure}
\hfill
\begin{subfigure}{0.49\textwidth}
\centering
\includegraphics[width=\textwidth]{step2_gmm_gamma_l_by_t.png}
\caption{Lead-time effect by auction date}
\end{subfigure}
\caption{Linear projection of participants' risk aversion}
\label{fig:gamma_trends}
\begin{minipage}{0.94\textwidth}
\footnotesize
\textit{Notes:} Equal-auction projection. Shaded regions are pointwise central
95 percent intervals from 200 global firm-pairs bootstrap draws.
\end{minipage}
\end{figure}

\begin{figure}[H]
\centering
\begin{subfigure}{0.49\textwidth}
\centering
\includegraphics[width=\textwidth]{step2_gmm_cost_t_l0.png}
\caption{Auction-date trend at zero lead time}
\end{subfigure}
\hfill
\begin{subfigure}{0.49\textwidth}
\centering
\includegraphics[width=\textwidth]{step2_gmm_cost_l_by_t.png}
\caption{Lead-time effect by auction date}
\end{subfigure}
\caption{Linear projection of participants' costs}
\label{fig:cost_trends}
\begin{minipage}{0.94\textwidth}
\footnotesize
\textit{Notes:} Equal-auction projection. The undiscounted specification in
Panel (b) divides cost by $\delta^{l_a}$ before projection. Shaded regions are
pointwise central 95 percent intervals from 200 global firm-pairs bootstrap
draws.
\end{minipage}
\end{figure}

The fitted gamma trend declines modestly from 2010 and is nearly flat over
2012--2015, but its confidence band is broad. The lead-time effect on gamma is
close to zero throughout. In contrast, the cost projection retains the marked
U-shape in auction date and the inverse-U lead-time pattern found in the paper;
both are much more precisely estimated.

\subsection{Step 2 Bootstrap Diagnostics}

The Step 2 bootstrap reproduces the Step 1 global firm-pairs draw. Firm
multiplicities enter the bid-price likelihood and concomitant models as
frequency weights and also enter the winner summaries. A draw can assign zero
weight to all observed winners in an auction. Such an auction is inactive in
that draw, and its nuisance parameters, cost inversions, and distributional
projection observations are omitted rather than replaced by original-sample
estimates.

All 200 bid-distribution problems converge locally. For the censored
prior-capacity model, 87 draws reach Ipopt's preferred local tolerance and 113
reach its acceptable tolerance; one draw uses the deterministic retry. No
conditional variance reaches its computational bound. The number of active
sealed-bid auctions ranges from 12 to 14. In 121 draws all 14 auctions are
active, 74 draws contain 13, and five draws contain 12. The corresponding
draws recover 549, 546, and 536 winners, respectively. This variation is the
intended conditional-on-auctions
firm-resampling design. Every aggregate statistic in Table
\ref{tab:participant_results} is estimable in all 200 draws.

\subsection{Sensitivity Analysis}

All 18 Step 2 sensitivity specifications recover all 549 eligible winners.
The main patterns are summarized in Table \ref{tab:sensitivity}; ranges are
taken over each stated grid.

\begin{table}[H]
\centering
\small
\caption{Step 2 sensitivity ranges for capacity-weighted winners}
\label{tab:sensitivity}
\begin{tabular}{lrrrr}
\toprule
Perturbation & Grid & Mean $\gamma$ & Risk premium & Mean cost\\
\midrule
Spot-price variance multiplier & 0.5--1.5 & 0.247--2.342 & 8.637--9.086 & 23.742--24.399\\
Spot-price mean shift & $-2$--$2$ SD & 0.457--0.617 & 7.091--9.577 & 23.941--24.009\\
Winning-bid variance multiplier & 0.5--1.5 & 0.554--0.566 & 8.597--8.790 & 23.985--23.989\\
Discount factor $\delta$ & 0.90--0.99 & 0.332--0.891 & 5.661--12.416 & 14.066--38.319\\
\bottomrule
\end{tabular}
\end{table}

Changes in assumed spot-price variance are largely offset by changes in
estimated gamma, leaving the economically relevant risk premium and recovered
cost comparatively stable. The winning-bid variance has little effect. For the
discount-factor comparison, Table \ref{tab:step2_delta_scaling} places the risk
premium in comparable annualized price units.

\begin{table}[H]
\centering
\small
\caption{Step 2 discount-factor sensitivity in comparable units}
\label{tab:step2_delta_scaling}
\begin{tabular}{rrrrrr}
\toprule
$\delta$ & Mean $\widetilde\delta_a$ & Risk premium
& Risk premium / $\widetilde\delta_a$ & Mean cost & Premium / cost\\
\midrule
0.90 & 0.320 & 5.661 & 17.528 & 14.066 & 0.402\\
0.95 & 0.550 & 8.711 & 15.817 & 23.987 & 0.363\\
0.99 & 0.883 & 12.416 & 14.057 & 38.319 & 0.324\\
\bottomrule
\end{tabular}
\begin{minipage}{0.92\textwidth}
\footnotesize
\textit{Notes:} Capacity-weighted winner statistics. The normalized premium
divides each winner's risk premium by its auction-specific
$\widetilde\delta_a$ before averaging.
\end{minipage}
\end{table}

The raw premium changes by a factor of 2.19 across the endpoint discount
factors, but the normalized premium changes by only a factor of 1.25. The
premium-to-cost ratio also moves much less than the component levels. Thus most
of the apparent sensitivity is the expected NPV scaling, rather than an
economically implausible change in the underlying price of risk. The remaining
movement reflects the nonlinear effect of $\delta$ on wholesale-price variance
and re-estimation of gamma.

\section{Reproducibility Files}

\begin{itemize}
\item \texttt{code/data\_processing/estimation\_prep.R}: estimation data and
prior started wind capacity;

\item \texttt{code/estimation/step1.jl}: original-sample and bootstrap GMM;

\item \texttt{code/estimation/step1\_result.jl}: bootstrap standard errors and
implied risk-aversion statistics;

\item \texttt{code/estimation/step1\_sa.jl}: Step 1 sensitivity estimates;

\item \texttt{code/estimation/step2\_dist.jl} and
\texttt{step2\_dist\_func.jl}: bid-price, capacity, and prior-capacity
concomitant distributions;

\item \texttt{code/estimation/step2.jl} and \texttt{step2\_func.jl}: winning
probabilities, cost inversion, and original/bootstrap outputs;

\item \texttt{code/estimation/step2\_sa.jl}: Step 2 sensitivity outputs;

\item \texttt{code/estimation/step2\_result.jl}: bootstrap standard errors and
aggregate winner and participant summaries;

\item \texttt{notes/gmm/step1\_gmm\_reduced\_form.R}: reduced form;

\item \texttt{notes/gmm/step1\_gmm\_gof.jl} and
\texttt{step1\_gmm\_gof.R}: predictive diagnostics;

\item \texttt{notes/gmm/step2\_gmm\_results.R}: participant-distribution
summaries, trend figures, and discount-factor normalization;

\item \texttt{notes/gmm/step2\_gmm\_*.csv} and
\texttt{step2\_gmm\_*.png}: note-specific Step 2 tables and figures;

\item \texttt{output/result/params1\_hat.csv},
\texttt{params1\_result.csv}, and \texttt{params1\_status.csv}: Step 1
estimates, bootstrap inference, and optimization statuses;

\item \texttt{output/result/params2\_dist\_result.csv} and
\texttt{params2\_dist\_status.csv}: Step 2 nuisance distributions;

\item \texttt{output/result/gamma\_c\_winner\_result.csv}: winner recovery;

\item \texttt{output/result/gamma\_c\_summary\_result.csv} and
\texttt{gamma\_c\_lp\_result.csv}: participant summaries and projections;

\item \texttt{output/result/step2\_status.csv}: Step 2 recovery diagnostics.
\end{itemize}

\end{document}
